\section{Обзор существующих решений}
\subsection{Анализаторы сетевого трафика}
Современные системы анализа сетевого трафика можно классифицировать по глубине инспектирования, архитектуре обработки и модели расширения функциональности. В данном подразделе рассматриваются три наиболее распространённых открытых решения: Ntopng, Suricata и Zeek, с акцентом на их подходы к генерации и обработке событий.

\subsubsection{Ntopng}
Ntopng представляет собой систему мониторинга сетевого трафика, ориентированную на визуализацию потоков, анализ пропускной способности и выявление аномалий на основе статистических моделей. Архитектура Ntopng базируется на потоковом анализе (flow-based), где пакеты агрегируются в записи NetFlow/IPFIX с последующим хранением в Redis или базе данных TimescaleDB. Генерация событий в Ntopng реализована преимущественно через пороговые механизмы и интеграцию с внешними системами оповещения. Расширяемость обеспечивается через плагины на языке Lua и REST API, однако внутренняя событийная шина не выделена в отдельный компонент, что ограничивает гибкость маршрутизации событий и требует перезапуска процессов при глубокой модификации логики обработки. Система эффективна для задач мониторинга производительности, но не предназначена для детального анализа протоколов прикладного уровня в реальном времени \cite{ntopng2023doc}.

\subsubsection{Suricata}
Suricata является многопоточным IDS/IPS-движком, поддерживающим сигнатурный анализ, обнаружение аномалий и глубокую инспекцию пакетов. Архитектура Suricata построена на основе конвейера обработки, где каждый пакет проходит через стадии дефрагментации, реконструкции потоков и применения правил. События генерируются при совпадении сигнатур или нарушении правил безопасности и выводятся в форматах JSON, Eve или syslog. Расширение функциональности возможно через добавление новых правил в формате YAML/Lua и подключение внешних модулей, однако базовый механизм событийного движка жёстко привязан к ядру обработки. Динамическое добавление новых типов событий требует компиляции или перезапуска, а обработка событий осуществляется синхронно в рамках потока обработки пакетов, что создаёт узкие места при высокой нагрузке. Suricata демонстрирует высокую пропускную способность, но её событийная модель не оптимизирована для гибкой аналитической кастомизации без вмешательства в исходный код \cite{suricata2023arch}.

\subsubsection{Zeek}
Zeek (ранее Bro) позиционируется как система мониторинга сетевой безопасности, ориентированная на протокольный анализ и генерацию структурированных логов. В отличие от сигнатурных систем, Zeek использует событийно-ориентированную модель, где каждый этап обработки пакета (установление соединения, запрос/ответ протокола, завершение сессии) генерирует соответствующее событие. Обработка событий осуществляется через встроенный скриптовый язык ZeekScript, позволяющий описывать логику анализа без модификации ядра. Однако интерпретация скриптов вносит значительные накладные расходы, а управление состоянием событий ограничено однопоточной или слабо масштабируемой моделью. Для внесения изменений в схему событий или добавления новых обработчиков часто требуется перезагрузка политик или перезапуск процесса, что снижает доступность в режимах непрерывного мониторинга. Zeek остаётся эталоном в области протокольной диссекции, но её событийная подсистема не удовлетворяет требованиям к высокой пропускной способности и динамической расширяемости в современных распределённых средах \cite{zeek2022design}.

\subsubsection{Итоговый анализ анализаторов}
Сравнительный анализ показывает, что существующие системы либо фокусируются на потоковой статистике (Ntopng), либо на сигнатурном обнаружении (Suricata), либо на скриптовой событийной обработке (Zeek). Ни одно из решений не предоставляет выделенной, высокопроизводительной и типобезопасной библиотеки событий, предназначенной исключительно для встраивания в сторонние DPI-конвейеры. Ограничения включают отсутствие динамического реестра схем событий, синхронную обработку с блокирующими вызовами, необходимость перезапуска для обновления логики и слабую изоляцию пользовательских обработчиков. Эти недостатки формируют технологический пробел, который заполняется разработкой независимой расширяемой системы генерации событий.

% \subsection{Библиотеки систем событий}
\subsection{Boost.Signals2}
Boost.Signals2 представляет собой библиотеку для реализации паттерна ``наблюдатель'' с поддержкой многопоточности и управления временем жизни подключений \cite{boost2023}. Библиотека обеспечивает типобезопасное соединение сигналов и слотов через механизм function objects.

Ключевыми преимуществами Boost.Signals2 являются зрелость реализации, кроссплатформенность и интеграция с экосистемой Boost. Однако библиотека демонстрирует значительные накладные расходы на диспетчеризацию (~50-100 нс на вызов обработчика) из-за использования \texttt{std::function} и динамического выделения памяти для хранения подключений \cite{schaling2011boost}.

Для сценариев DPI с требованиями к задержке обработки на уровне единиц наносекунд накладные расходы Boost.Signals2 являются неприемлемыми, особенно при большом числе обработчиков на событие.

\subsection{Eventpp}
Eventpp — современная библиотека обработки событий для C++, ориентированная на производительность и типобезопасность \cite{eventpp2023}. Библиотека использует шаблоны для генерации кода диспетчеризации на этапе компиляции, что позволяет избежать накладных расходов type erasure.

Eventpp поддерживает различные стратегии выполнения обработчиков, включая синхронный вызов и асинхронную очередь. Однако библиотека не предоставляет встроенной интеграции с асинхронными runtime (asio, libuv), что ограничивает её применимость для гибридных сценариев обработки.

Кроме того, шаблонная природа Eventpp приводит к увеличению времени компиляции и размера бинарного файла при большом числе типов событий, что может быть критичным для встраиваемых систем.

\subsection{Libuv}
Libuv представляет собой кроссплатформенную библиотеку асинхронного ввода-вывода, используемую в проекте Node.js \cite{libuv2023}. Библиотека предоставляет абстракции над системными механизмами асинхронности (epoll, kqueue, IOCP) и реализует event loop с поддержкой таймеров, файловых дескрипторов и межпоточной коммуникации.

Ключевым преимуществом Libuv является высокая производительность асинхронных операций и зрелая реализация event loop. Однако библиотека ориентирована преимущественно на IO-bound сценарии и не предоставляет оптимальных механизмов для синхронной обработки CPU-bound задач в контексте вызова.

Для гибридных сценариев DPI, требующих как детерминированной обработки критичных по времени операций, так и асинхронного выполнения длительных задач, использование Libuv в качестве единственной основы системы событий является неоптимальным.

\subsection{Итоговый анализ библиотек событий}
Анализ существующих библиотек обработки событий позволяет сформулировать требования к целевому решению:

\begin{table}[H]
\centering
\caption{Сравнительный анализ библиотек событий}
\label{tab:libs_comparison}
\begin{tabularx}{\textwidth}{lXXX}
\toprule
\textbf{Критерий} & \textbf{Boost.Signals2} & \textbf{Eventpp} & \textbf{Libuv} \\
\midrule
Накладные расходы на вызов & 50-100 нс & 5-15 нс & 10-30 нс \\
Поддержка async runtime & Ограниченная & Отсутствует & Полная \\
Типобезопасность & Высокая & Очень высокая & Средняя \\
Гибкость выполнения & Низкая & Средняя & Высокая \\
Интеграция с DPI & Низкая & Средняя & Средняя \\
\bottomrule
\end{tabularx}
\end{table}

Ни одна из рассмотренных библиотек не удовлетворяет одновременно всем требованиям современных DPI-систем: минимальные накладные расходы, гибридная модель выполнения, типобезопасность и гибкость расширения. Это обосновывает необходимость разработки специализированного решения.

\newpage